% source: erdos-872/current_state.md (Polynomial shield lower bound --- Theorem A)
% source: erdos-872/researcher-13-pro-upper-half-fan-lower-bound.md
% source: erdos-872/phase4/t2_core_paper_note.md
% source: erdos-872/lean/README.md

\section{Lower Bounds}\label{sec:lower-bounds}

\subsection{The polynomial shield barrier and T1}

The shield reduction of \Cref{thm:shield-reduction} turns lower bounds on
$\L(n)$ into upper bounds on the weighted lower-half parameter $\beta(P)$.
Theorem~A shows that this route already faces a sharp obstruction at polynomial
shield size.

\begin{theorem}[Theorem A]\label{thm:theorem-a}
Fix $0<\alpha<1$.  For every $P \subseteq \U$ with $|P| \le n^\alpha$,
\[
  \beta(P) \ge \left(\frac12 \log\frac1\alpha + o(1)\right)n.
\]
\end{theorem}

\begin{proof}
Fix $\delta$ with $\alpha<\delta<1$, and let
\[
  Q_\delta(P)
  :=
  \{p \le n^\delta : p \text{ prime and } p \nmid u \text{ for all } u \in P\}
\]
be the set of small primes not covered by $P$.  Since distinct primes are
incomparable, $Q_\delta(P)$ is an antichain in $L(P)$.  Therefore
\[
  \beta(P) \ge \sum_{p \in Q_\delta(P)} w_n(p).
\]
For a prime $p \le n^\delta$ we have
\[
  w_n(p)=\#\{u \in (n/2,n] : p \mid u\}
  = \frac{n}{2p}+O(1),
\]
uniformly in $p$.  Hence
\[
  \beta(P)
  \ge
  \frac n2 \sum_{p \in Q_\delta(P)} \frac1p + O(\pi(n^\delta)).
\]
Since $\pi(n^\delta)=o(n)$, it suffices to lower-bound the reciprocal mass of
$Q_\delta(P)$.

Let $C_\delta(P)$ be the complementary set of covered primes
\[
  C_\delta(P)
  :=
  \{p \le n^\delta : p \text{ prime and } p \mid u \text{ for some } u \in P\}.
\]
Each covered prime $p$ is assigned to one $u \in P$ divisible by $p$.  Because
the assigned primes are distinct and all divide $u$, their product divides $u$,
so the sum of their logarithms is at most $\log u \le \log n$.  Summing over
$u \in P$ gives the global log-budget inequality
\[
  \sum_{p \in C_\delta(P)} \log p \le |P| \log n \le n^\alpha \log n.
\]

Among all prime sets with a fixed logarithmic budget, the reciprocal mass
$\sum 1/p$ is maximized by taking the smallest primes first.  Therefore the
covered set with largest possible reciprocal mass is the initial prime segment
up to about $n^\alpha$, and Mertens' theorem yields
\[
  \sum_{p \in C_\delta(P)} \frac1p \le \log\log(n^\alpha) + O(1)
  = \log\log n + \log \alpha + O(1).
\]
On the other hand,
\[
  \sum_{p \le n^\delta} \frac1p = \log\log(n^\delta)+O(1)
  = \log\log n + \log\delta + O(1).
\]
Subtracting gives
\[
  \sum_{p \in Q_\delta(P)} \frac1p
  \ge \log\frac{\delta}{\alpha} + o(1).
\]
Substituting into the lower bound for $\beta(P)$,
\[
  \beta(P) \ge \left(\frac12 \log\frac{\delta}{\alpha}+o(1)\right)n.
\]
Now let $\delta \uparrow 1$ to obtain
\[
  \beta(P) \ge \left(\frac12 \log\frac1\alpha+o(1)\right)n.
\]
\end{proof}

\begin{corollary}[T1]\label{cor:t1}
There exists $c>0$ such that
\[
  \L(n) \ge c \frac{n\log\log n}{\log n}.
\]
More precisely,
\[
  \L(n) \ge \left(\frac18-o(1)\right)\frac{n\log\log n}{\log n}.
\]
\end{corollary}

\begin{proof}
Choose a shield family $P$ of size on the order of $\log n$.  Then
$|P|=n^{o(1)}$, so Theorem~\ref{thm:theorem-a} gives
\[
  \beta(P) \ge \left(\frac12+o(1)\right)n\log\log n.
\]
The calibrated shield construction from the current-state note converts this
weighted lower-half mass into a game-length lower bound of size
$(1/8-o(1))n\log\log n/\log n$ via \Cref{thm:shield-reduction}.  The
substantive point is that the shield-reduction architecture already excludes the
possibility $\L(n)=O(n/\log n)$.
\end{proof}

\begin{remark}
Theorem~\ref{thm:theorem-a} is strategy-independent: it is a static theorem
about weighted primitive sets.  The dependence on the actual game enters only
when a specific Prolonger shield family is fed into
\Cref{thm:shield-reduction}.  In particular, the exponent $1/e$ barrier for
linear lower bounds is a property of the reduction itself, not of any single
failed strategy.
\end{remark}

\subsection{T2: a Maker-first hypergraph capture lower bound}

The T2 lower bound refines T1 by moving from one-dimensional shield accounting
to a two-stage game.  First Prolonger activates many useful small-prime pairs
$(a,c)$.  Then each activated pair carries a fiber of upper-half residual
targets
\[
  t = acb \in \left(\frac n2,n\right],
\]
where $b$ is a large prime.  The central point is that the residual game on
these targets is no harder for Prolonger than a scored Maker-first
3-uniform-hypergraph game.

We record the finite abstract lemmas exactly as they are used.

\begin{lemma}[Weighted Maker-first pair capture]\label{lem:t2-graph}
In the finite weighted graph game recorded in the consolidated T2 note, Maker
has a strategy that secures at least one eighth of the initial total
edge-weight.
\end{lemma}

\begin{lemma}[Weighted Maker-first scored 3-uniform capture]\label{lem:t2-hyper}
In the corresponding finite scored 3-uniform hypergraph game, Maker has a
strategy that secures at least one eighth of the initial total hyperedge-weight.
\end{lemma}

The residual divisibility game is compared to the hypergraph model through the
following theorem.  Its local arithmetic and family-level comparison layer are
substantially formalized in Lean.

\begin{theorem}[Residual comparison theorem]\label{thm:t2-residual}
Let $\mathcal T_*$ be a residual family of targets $t=acb \in (n/2,n]$ with
$a<c\le Y=n^\delta$ prime and $b>Y$ prime, arranged so that $a$, $c$, and $ac$
are already unavailable from the activation stage.  Let $H_*$ be the
3-uniform hypergraph with one hyperedge
\[
  e_t = \{b,ab,cb\}
\]
for each $t \in \mathcal T_*$.  Then:
\begin{enumerate}
  \item the only harmful future moves for $t$ are $b$, $ab$, $cb$, and $t$
    itself;
  \item the moves $b$, $ab$, and $cb$ delete exactly the incident hyperedges in
    $H_*$;
  \item an exact target move $t$ hits only its own edge;
  \item every live hyperedge in $H_*$ corresponds to a legal actual Prolonger
    move in the divisibility game.
\end{enumerate}
Consequently the actual residual divisibility game is Maker-friendlier than the
scored hypergraph game on $H_*$.
\end{theorem}

\begin{proof}[Proof sketch]
Fix $t=acb \in \mathcal T_*$.  Since $t>n/2$, it has no proper multiple in
$\{2,\dots,n\}$.  Hence any harmful move must be either $t$ itself or a proper
divisor of $t$.  The positive divisors of $t$ are
\[
  1,\ a,\ c,\ ac,\ b,\ ab,\ cb,\ t,
\]
and $1$ is not a legal move while $a,c,ac$ are already unavailable.  This gives
the harmful set $\{b,ab,cb,t\}$.

Next, a move $b$ kills exactly those targets with the same large prime $b$,
because $b>Y\ge a,c$ cannot divide the small part of another target.  The same
argument shows that $ab$ and $cb$ kill exactly the targets using those slots.
Thus the slot moves are exactly vertex deletions in the hypergraph model.

Distinct targets in $(n/2,n]$ are incomparable, so an exact target move cannot
kill any other target.  Finally, if the hyperedge $e_t$ is live, none of
$b,ab,cb,t$ has been played, while $a,c,ac$ were unavailable from the start.
There is therefore no played number comparable with $t$, so $t$ is legal.  This
proves that every hypergraph play transfers to the actual divisibility game and
that the hypergraph model is Maker-friendlier.
\end{proof}

The activation stage keeps track of how much residual target mass survives on
secured pairs.

\begin{theorem}[Activation-stage inequality]\label{thm:t2-activation}
Let
\[
  W_0 := \sum_{a<c\le Y} |J_{a,c}\cap\PP|
\]
be the initial total token weight, where $J_{a,c}$ is the large-prime interval
attached to the pair $(a,c)$.  Let $Q_{\mathrm{end}}$ be the graph-potential
quantity at the end of the activation stage, and let $E$ be the total number of
tokens deleted by off-model Shortener moves during activation.  Then
\[
  Q_{\mathrm{end}} \ge \frac{W_0}{8} - E.
\]
Moreover, for every fixed $\delta<1/4$,
\[
  E = o\!\left(\frac{n(\log\log n)^2}{\log n}\right).
\]
\end{theorem}

\begin{proof}[Proof sketch]
Ignoring the off-model deletions, the activation game is exactly the weighted
Maker-first graph game of \Cref{lem:t2-graph}: when Prolonger claims an edge he
converts one token into score and leaves the remaining tokens on that edge, so
the same one-step domination argument keeps the graph potential nondecreasing.
This yields the clean lower bound $Q_{\mathrm{end}}\ge W_0/8$.

Restoring the off-model deletions costs at most one unit of potential per
deleted token, so
\[
  Q_{\mathrm{end}} \ge \frac{W_0}{8}-E.
\]
To estimate $E$, observe that there are at most
\[
  O\!\left(\frac{Y^2}{\log^2 Y}\right)
\]
activation rounds.  A large-prime deletion kills at most one token per pair,
namely
\[
  O\!\left(\frac{Y^2}{\log^2 Y}\right)
\]
tokens; a lateral move $pb$ kills only
$O(Y/\log Y)$ targets; and an exact target kills one token.  Hence
\[
  E \ll \frac{Y^4}{\log^4 Y}.
\]
With $Y=n^\delta$ and $\delta<1/4$, this is
\[
  o\!\left(\frac{n(\log\log n)^2}{\log n}\right).
\]
\end{proof}

\begin{theorem}[T2]\label{thm:t2}
For every fixed $0<\delta<1/4$ there exists $c_\delta>0$ such that
\[
  \L(n) \ge c_\delta \frac{n(\log\log n)^2}{\log n}
\]
for all sufficiently large $n$.
\end{theorem}

\begin{proof}
At the end of the activation stage let
\[
  M := \sum_{e\in C_{\mathrm{end}}} w_{\mathrm{end}}(e)
\]
be the residual live target weight on secured pairs.  By
\Cref{thm:t2-activation},
\[
  M \ge \frac{W_0}{8} - E - S_{\mathrm{end}},
\]
where $S_{\mathrm{end}}$ is the number of activation moves already converted
into score.  Since
\[
  S_{\mathrm{end}} = O\!\left(\frac{Y^2}{\log^2 Y}\right)
  = o\!\left(\frac{n(\log\log n)^2}{\log n}\right),
\]
and the number-theoretic part of the construction gives
\[
  W_0 \gg_\delta \frac{n(\log\log n)^2}{\log n},
\]
we conclude that
\[
  M \gg_\delta \frac{n(\log\log n)^2}{\log n}.
\]

Now transfer to the actual residual game using
\Cref{thm:t2-residual}.  The hypergraph game on the family of residual targets
is Maker-friendlier than the actual divisibility game, so Maker can secure at
least one eighth of the residual weight by \Cref{lem:t2-hyper}.  Therefore the
number of actual game moves remaining after activation is at least
\[
  \frac{M}{8} \gg_\delta \frac{n(\log\log n)^2}{\log n},
\]
and the same lower bound holds for $\L(n)$.
\end{proof}

\begin{remark}
The present status of T2 is slightly different from that of T1.  The finite
graph and hypergraph potential cores, the local arithmetic of the divisibility
embedding, and the family-level residual comparison layer are Lean-backed.  The
remaining unformalized part is the activation bookkeeping and the final
packaging of the Maker-first wrapper into one theorem.  For the purposes of the
paper, the proof chain is rigorous prose with substantial formal support rather
than a fully closed Lean artifact.
\end{remark}
